\section{Дополнительные сведения об экспериментах}\label{sec:appendix}
\subsection{Подготовка данных и параметры}
Источники наборов данных и модель текстовых представлений описаны в~\citep{source_agnews,source_banking,source_sklearn,source_drybean,source_minilm}. Подвыборки отбираются без возвращения с начальным значением 2026 после очистки. Из числовых данных удаляются совпадающие строки. При подготовке эталонных данных со всеми метками конфликтующие совпадения исключаются, затем метки скрываются. Правило выбора получает только раскрытую разметку.

Для AG News и Banking77 до очистки объединяются исходные обучающая и тестовая части. Из 20 Newsgroups удаляются заголовки, подписи и цитаты. Нормализуются пробелы, исключаются пустые и конфликтующие тексты, точные дубликаты объединяются. Известные связи документов учитываются при разделении. Во всех отобранных подвыборках группы оказались одноэлементными.

Стандартизация числовых признаков обучается на всём наборе без меток. Для текстов используются нормированные представления all-MiniLM-L6-v2, ограничение длины 256 токенов и размер пакета 64. PCA использует полный SVD без приведения дисперсий компонент к единице. После проекции на 50 компонент повторная нормировка отсутствует.

KMeans выполняет 20 инициализаций и до 300 итераций. Спектральная кластеризация использует граф 15 ближайших соседей и 10 инициализаций заключительного KMeans. Начальное значение генератора алгоритмов равно 42. Разделение резерва и проверочной части выполняется с начальным значением 31415 и целевой долей проверочных групп $0{,}5$. Фактические размеры показаны в табл.~\ref{tab:split}.

\begin{longtable}{lrrrrr}
\caption{Фактические группы и размеры разделения.}\label{tab:split}\\
\toprule
Набор & $n$ & Групп & Макс. размер & $|S|$ & $|T|$\\
\midrule\endfirsthead
\multicolumn{6}{l}{Таблица \thetable\ (продолжение)}\\

Набор & $n$ & Групп & Макс. размер & $|S|$ & $|T|$\\\midrule\endhead
Breast Cancer & 569 & 569 & 1 & 284 & 285\\
Digits & 600 & 600 & 1 & 300 & 300\\
Dry Bean & 600 & 600 & 1 & 300 & 300\\
AG News & 600 & 600 & 1 & 300 & 300\\
20 Newsgroups & 600 & 600 & 1 & 300 & 300\\
Banking77 & 600 & 600 & 1 & 300 & 300\\
\bottomrule\end{longtable}

Все раскрытия основной серии выбираются из резерва без возвращения. Бюджеты 20 и 50 вложены в пределах повтора. Бутстреп перевыбирает раскрытые объекты с возвращением. Для сравнений кандидатов используются общие кратности. Во всех 24000 перевыборках основной серии присутствуют как минимум два класса и одноимённые пары. Две исходные разметки Banking77 при 20 метках состоят из одноэлементных классов. Повторения объектов при бутстрепе создают положительные кратности пар, сохраняя исходную информацию.

\subsection{Неопределённость среднего эффекта}
Интервалы среднего получены 10000 перевыборками двадцати парных эффектов с возвращением, начальное значение 20261112. Уровень 95\%, границы --- процентили 2,5 и 97,5. Данные, кандидаты и проверочная часть фиксированы. Интервалы вероятности ухудшения рассчитаны методом Клоппера---Пирсона отдельно на всех раскрытиях и при переключении. Поправка на множественные сравнения отсутствует, интервалы используются описательно.

\begingroup\footnotesize\setlength{\tabcolsep}{3pt}
\begin{longtable}{lrrrrrrr}
\caption{Условные интервалы среднего эффекта и верхние границы вероятностей ухудшения.}\label{tab:uncertainty}\\
\toprule
Набор & $m$ & $G_{.025}$ & $G_{.975}$ & Ухудш. & $p_{\max}$ & При смене & $p^s_{\max}$\\
\midrule\endfirsthead
\multicolumn{8}{l}{Таблица \thetable\ (продолжение)}\\

Набор & $m$ & $G_{.025}$ & $G_{.975}$ & Ухудш. & $p_{\max}$ & При смене & $p^s_{\max}$\\\midrule\endhead
20 Newsgroups & 20 & $0{,}0000$ & $0{,}0271$ & 0/20 & $0{,}1684$ & 0/2 & $0{,}8419$\\
20 Newsgroups & 50 & $0{,}0000$ & $0{,}0328$ & 0/20 & $0{,}1684$ & 0/3 & $0{,}7076$\\
AG News & 20 & $-0{,}0237$ & $0{,}0548$ & 1/20 & $0{,}2487$ & 1/3 & $0{,}9057$\\
AG News & 50 & $0{,}0000$ & $0{,}0303$ & 0/20 & $0{,}1684$ & 0/1 & $0{,}9750$\\
Banking77 & 20 & $0{,}0075$ & $0{,}0558$ & 0/20 & $0{,}1684$ & 0/5 & $0{,}5218$\\
Banking77 & 50 & $0{,}0000$ & $0{,}0000$ & 0/20 & $0{,}1684$ & --- & ---\\
Breast Cancer & 20 & $0{,}0000$ & $0{,}0216$ & 0/20 & $0{,}1684$ & 0/3 & $0{,}7076$\\
Breast Cancer & 50 & $0{,}0000$ & $0{,}0108$ & 0/20 & $0{,}1684$ & 0/1 & $0{,}9750$\\
Digits & 20 & $0{,}0049$ & $0{,}0535$ & 0/20 & $0{,}1684$ & 0/6 & $0{,}4593$\\
Digits & 50 & $0{,}0037$ & $0{,}0134$ & 0/20 & $0{,}1684$ & 0/7 & $0{,}4096$\\
Dry Bean & 20 & $-0{,}1063$ & $-0{,}0360$ & 12/20 & $0{,}8088$ & 12/12 & $1{,}0000$\\
Dry Bean & 50 & $-0{,}0269$ & $-0{,}0016$ & 5/20 & $0{,}4910$ & 5/5 & $1{,}0000$\\
\bottomrule\end{longtable}
\endgroup
Если все двадцать наблюдённых эффектов нулевые, эмпирический интервал среднего равен $[0;0]$. За пределами наблюдений возможны редкие исходы и ненулевой средний эффект. При нуле ухудшений в 20 раскрытиях верхняя граница вероятности равна $0{,}1684$. Для единственного переключения без ухудшения условная верхняя граница достигает $0{,}975$.

\subsection{Составляющие потери и состав разметки}
Разложение~\eqref{eq:regret-decomposition} вычисляется для каждого раскрытия. Средние компоненты показаны в табл.~\ref{tab:main-loss}. Они характеризуют ошибку допуска и выбор внутри него при одном наборе кандидатов.

\begingroup\small
\begin{exampletable}\small
\caption{Составляющие потери M6 в основной серии.}\label{tab:main-loss}
\begin{tabular}{lrrr}
\toprule Набор & $m$ & $D_{\mathrm{exclude}}$ & $D_{\mathrm{inside}}$\\\midrule
20 Newsgroups & 20 & $0{,}0164$ & $0{,}0000$\\
20 Newsgroups & 50 & $0{,}0000$ & $0{,}0000$\\
AG News & 20 & $0{,}0202$ & $0{,}0214$\\
AG News & 50 & $0{,}0000$ & $0{,}0000$\\
Banking77 & 20 & $0{,}0000$ & $0{,}0000$\\
Banking77 & 50 & $0{,}0000$ & $0{,}0000$\\
Breast Cancer & 20 & $0{,}0000$ & $0{,}0000$\\
Breast Cancer & 50 & $0{,}0000$ & $0{,}0000$\\
Digits & 20 & $0{,}0037$ & $0{,}0000$\\
Digits & 50 & $0{,}0000$ & $0{,}0000$\\
Dry Bean & 20 & $0{,}0049$ & $0{,}0826$\\
Dry Bean & 50 & $0{,}0057$ & $0{,}0147$\\
\bottomrule\end{tabular}\end{exampletable}
\endgroup

В проверке состава разметки при $m=20$ случайно выбираются 2, 4 или 6 классов с квотами $(10,10)$, $(5,5,5,5)$ и $(4,4,3,3,3,3)$. Это даёт 90, 40 и 24 одноимённые пары. Для всех квот допускаются категории с не менее чем десятью объектами в резерве. На Digits и Dry Bean этому условию удовлетворяют все классы. Для каждого условия выполнены 20 раскрытий. При каждом повторе категории выбираются заново равномерно без возвращения.

\begingroup\small
\begin{exampletable}\small
\caption{Состав разметки при m=20. По 20 раскрытий, условные 95\%-интервалы среднего.}\label{tab:coverage-control}
\begin{tabular}{lrrrrrr}
\toprule Набор & Классов & $G$ & Низ & Верх & Смен & Ухудш.\\\midrule

Digits & 2 & $0{,}0154$ & $-0{,}0495$ & $0{,}0762$ & 6 & 2\\
Digits & 4 & $0{,}0039$ & $-0{,}0255$ & $0{,}0390$ & 6 & 2\\
Digits & 6 & $-0{,}0194$ & $-0{,}0527$ & $0{,}0140$ & 11 & 5\\
Dry Bean & 2 & $-0{,}0096$ & $-0{,}0416$ & $0{,}0174$ & 9 & 3\\
Dry Bean & 4 & $-0{,}0772$ & $-0{,}1227$ & $-0{,}0360$ & 12 & 11\\
Dry Bean & 6 & $-0{,}0023$ & $-0{,}0392$ & $0{,}0338$ & 9 & 6\\
\bottomrule\end{tabular}
\end{exampletable}
\endgroup
Скрытые метки используются только для формирования условий раскрытия. Правило получает раскрытые значения. Вместе с покрытием меняются квоты, число пар и смысл категорий. Самостоятельный эффект покрытия по этому опыту не установлен. Все дополнительные проверки выполнены после анализа основной серии.

\subsection{Вычислительная устойчивость бутстрепа}
Использованы все 80 разметок Digits и Breast Cancer: по 20 раскрытий и два бюджета. Для каждой выполнены расчёты с $B\in\{30,100,300\}$ и начальными значениями 0--4. Ориентир использует $B=3000$ и начальное значение 99. При одинаковом начальном значении последовательности перевыборок для разных $B$ вложены.

\begingroup\small
\begin{longtable}{lrrrr}
\caption{Совпадения с вычислительным ориентиром при B=3000.}\label{tab:b-control}\\
\toprule
Набор & $m$ & $B$ & Совпадений & Запусков\\
\midrule\endfirsthead
\multicolumn{5}{l}{Таблица \thetable\ (продолжение)}\\

Набор & $m$ & $B$ & Совпадений & Запусков\\\midrule\endhead
Breast Cancer & 20 & 30 & 100 & 100\\
Breast Cancer & 20 & 100 & 100 & 100\\
Breast Cancer & 20 & 300 & 100 & 100\\
Breast Cancer & 50 & 30 & 100 & 100\\
Breast Cancer & 50 & 100 & 100 & 100\\
Breast Cancer & 50 & 300 & 100 & 100\\
Digits & 20 & 30 & 96 & 100\\
Digits & 20 & 100 & 96 & 100\\
Digits & 20 & 300 & 96 & 100\\
Digits & 50 & 30 & 100 & 100\\
Digits & 50 & 100 & 100 & 100\\
Digits & 50 & 300 & 100 & 100\\
\bottomrule\end{longtable}
\endgroup
На Digits при 20 метках все несовпадения относятся к повторам 7 и 11. При несовпадении ARI ниже ориентира на $0{,}0244$ и $0{,}1579$ соответственно. Одинаковая доля совпадений поэтому скрывает разные размеры потери. В остальных 78 разметках совпадение полное.

Ориентир дополнительно рассчитан при начальных значениях 99--103. В 79 разметках все пять решений совпали. В повторе 7 Digits при 20 метках получены два решения: кандидат находится около границы допуска. Рост числа перевыборок повышает точность вычисления условной SD, однако устойчивость окончательного выбора в рассмотренном диапазоне остаётся отдельным свойством.

\subsection{Конечный пример для H3}
Рассмотрим шесть объектов с метками $(0,0,1,1,2,2)$. Резерв $S=\{1,2,3,4\}$ полностью размечен, $L=S$, проверочная часть $T=\{5,6\}$. Кандидаты имеют номера кластеров $(0,0,1,1,2,2)$ и $(0,0,1,1,0,1)$. На $L$ частичный ARI обоих равен 1, разность равна нулю во всех объектных перевыборках. На $T$ их ARI равны 1 и 0. Пример устанавливает возможность разного скрытого качества при нулевом разбросе. Внутренние оценки здесь не заданы, поэтому частота ошибочного окончательного выбора требует отдельного исследования.
